% Part: incompleteness
% Chapter: representability-in-q
% Section: sigma1-completeness

\documentclass[../../../include/open-logic-section]{subfiles}

\begin{document}

\olfileid{inc}{inp}{s1c}
\olsection{\texorpdfstring{$\Sigma_1$}{Sigma-1} પૂર્ણતા}

$\Th{Q}$ અને તેના સુસંગત,
સ્વયંસિદ્ધિગત વિસ્તરણો અપૂર્ણ
હોવા છતાં, આપણે જોયું છે કે
$\Th{Q}$ અંકો વિશેનાં ઘણાં
મૂળભૂત તથ્યો સાબિત કરે છે.
હકીકતમાં આ વાતને ખાસ્સી
આગળ લઈ જઈ શકાય છે.
$\Th{Q}$માં શું સાબિત કરી
શકાય છે તેનો વ્યાપ સમજવા
$\Delta_0$, $\Sigma_1$ અને
$\Pi_1$ !!{formula}sની
વિભાવનાઓ રજૂ કરીએ.
સ્થૂળ રીતે કહીએ તો
$\Sigma_1$ !!{formula} એ
$\lexists[x][!B(x)]$ આકારનું
છે, જ્યાં~$!B$ માત્ર વિધાનાત્મક
સંયોજકો અને મર્યાદિત પરિમાણકો
વડે રચાયેલું છે. આપણે બતાવીશું
કે $!A$ એ $\Struct{N}$માં
સાચું $\Sigma_1$ !!{sentence}
હોય, તો $\Th{Q} \Proves !A$
(\olref{thm:sigma1-completeness}).

\begin{defn}
\ollabel{defn:bd-quant}
\emph{મર્યાદિત અસ્તિત્વવાચક
!!{formula}} એ
$\lexists[x][(x < t \land !A(x))]$
આકારનું છે, જ્યાં~$t$ કોઈપણ
પદ છે; પરંપરાગત રીતે તેને
$\bexists{x < t}{!A(x)}$
લખીએ છીએ.
%
\emph{મર્યાદિત સર્વવાચક
!!{formula}} એ
$\lforall[x][(x < t \lif !A(x))]$
આકારનું છે, જ્યાં~$t$ કોઈપણ
પદ છે; પરંપરાગત રીતે તેને
$\bforall{x < t}{!A(x)}$
લખીએ છીએ.
\end{defn}

\begin{defn}
\ollabel{defn:delta0-sigma1-pi1-frm}
!!{formula}~$!B$ એ $\Delta_0$
ત્યારે કહેવાય જ્યારે તે
પરમાણુ !!{formula}sમાંથી
માત્ર વિધાનાત્મક સંયોજકો
અને મર્યાદિત પરિમાણીકરણ
વડે રચાયેલું હોય.
%
!!{formula}~$!A$ એ $\Sigma_1$
ત્યારે કહેવાય જ્યારે
$!A \ident \lexists[x][!B(x)]$
અને~$!B$ એ $\Delta_0$ હોય.
%
!!{formula}~$!A$ એ $\Pi_1$
ત્યારે કહેવાય જ્યારે
$!A \ident \lforall[x][!B(x)]$
અને~$!B$ એ $\Delta_0$ હોય.
\end{defn}

\begin{lem}
\ollabel{lem:q-proves-clterm-id}
માનો કે~$t$ બંધ પદ છે અને
$\Value{t}{N} = n$. તો
$\Th{Q} \Proves \eq[t][\num n]$.
\end{lem}

\begin{proof}
$t$ની જટિલતા પર અનુમાન
વડે આ સાબિત કરીએ.
આરંભિક કિસ્સામાં
$\Value{\Obj 0}{N} = 0$,
અને $\num 0 \ident \Obj 0$
હોવાથી
$\Th{Q} \Proves \eq[\Obj 0][\num 0]$.
%
અનુમાનના કિસ્સામાં માનો કે
$t_1$ અને~$t_2$ એવાં પદો છે કે
$\Value{t_1}{N} = n_1$,
$\Value{t_2}{N} = n_2$,
$\Th{Q} \Proves \eq[t_1][\num n_1]$,
અને $\Th{Q} \Proves
\eq[t_2][\num n_2]$.

તો $\Value{(t_1')}{N} = n_1 + 1$.
અનુમાનની ધારણા અને
!!{formula}~$\eq[\num{n_1}'][\num{n_1}']$
પર સમતા માટેના પ્રથમ-ક્રમના
નિયમો લગાડવાથી
$\Th{Q} \Proves
\eq[t_1'][{\num n_1}']$.
તેથી અંકોની વ્યાખ્યા પ્રમાણે
$\Th{Q} \Proves
\eq[t_1'][\num{n_1 + 1}]$.

સરવાળા માટે
\[
      \Value{(t_1 + t_2)}{N}
    = \Value{t_1}{N} + \Value{t_2}{N}
    = n_1 + n_2.
\]
અનુમાનની ધારણા તથા
સમતાના નિયમો વડે પહેલાં
$\Th{Q} \Proves
\eq[t_1 + t_2][\num{n_1} + t_2]$
અને તે નિયમો બીજી વાર
લગાડતાં
$\Th{Q} \Proves
\eq[t_1 + t_2][\num{n_1} + \num{n_2}]$
મળે છે.
\olref[inc][req][bre]{lem:q-proves-add}
પ્રમાણે
$\Th{Q} \Proves
\eq[\num{n_1} + \num{n_2}][\num{n_1 + n_2}]$;
તેથી $\Th{Q} \Proves
\eq[t_1 + t_2][\num{n_1 + n_2}]$.

$\times$ માટે પણ
\olref[inc][req][bre]{lem:q-proves-mult}
વાપરીને આવી જ દલીલ થાય છે.
%
અંકગણિતનાં બધાં બંધ પદો
આ કિસ્સાઓમાં આવી જાય છે;
તેથી $\Value{t}{N} = n$
હોય એવા દરેક બંધ પદ~$t$
માટે $\Th{Q} \Proves
\eq[t][\num n]$.
\end{proof}
 
\begin{prob}
$\times$ સંબંધિત
\olref{lem:q-proves-clterm-id}ની
સાબિતી વિગતે આપો.
\end{prob}

\begin{lem}
\ollabel{lem:atomic-completeness}
માનો કે $t_1$ અને~$t_2$
બંધ પદો છે. તો
\begin{enumerate}
\item જો $\Value{t_1}{N} = \Value{t_2}{N}$,
    તો $\Th{Q} \Proves \eq[t_1][t_2]$.
\item જો $\Value{t_1}{N} \neq \Value{t_2}{N}$,
    તો $\Th{Q} \Proves \eq/[t_1][t_2]$.
\item જો $\Value{t_1}{N} < \Value{t_2}{N}$,
    તો $\Th{Q} \Proves t_1 < t_2$.
\item જો $\Value{t_2}{N} \leq \Value{t_1}{N}$,
    તો $\Th{Q} \Proves \lnot(t_1 < t_2)$.
\end{enumerate}
\end{lem}

\begin{proof}
$t_1$ અને~$t_2$ આપેલાં હોય,
તો $n = \Value{t_1}{N}$ અને
$m = \Value{t_2}{N}$ નક્કી કરીએ.

માનો કે $!A \ident t_1 = t_2$.
\olref{lem:q-proves-clterm-id}
પ્રમાણે $\Th{Q} \Proves
\eq[t_1][\num n]$ અને
$\Th{Q} \Proves \eq[t_2][\num m]$.
જો $n = m$, તો
$\Th{Q} \Proves \eq[\num n][\num m]$;
અને સમતાની સંક્રામકતા વડે
$\Th{Q} \Proves \eq[t_1][t_2]$.
જો $n \neq m$, તો
$\Th{Q} \Proves \eq/[\num n][\num m]$;
સમતાની સંક્રામકતા ફરીથી
લગાડતાં $\Th{Q} \Proves
\eq/[t_1][t_2]$.
\sourcecorrection{OLINC-032}{સ્થિર અંગ્રેજી
મૂળ અહીં $t_2$નું મૂલ્ય~$m$
વ્યાખ્યાયિત કરે છે છતાં
$t_2=\num n$ લખે છે.
આ સમીકરણ $t_2=\num m$
કર્યું છે, જેથી $n=m$ અને
$n\neq m$ બંને કિસ્સા
યોગ્ય રીતે ચાલે.}

હવે $!A \ident t_1 < t_2$
લો. બંને કિસ્સામાં સ્વયંસિદ્ધિ
$!Q_8$ વાપરીશું; તે દરેક
$x,y$ માટે
$x < y \liff
\lexists[z][\eq[z' + x][y]]$
કહે છે.

માનો કે $\Sat{N}{t_1 < t_2}$.
તો કોઈ~$k \in \Nat$ માટે
$n + k + 1 = m$.
\olref{lem:q-proves-clterm-id}
પ્રમાણે $\Th{Q} \Proves
\eq[t_1][\num n]$ અને
$\Th{Q} \Proves \eq[t_2][\num m]$.
આ લેમ્માના પહેલા ભાગથી
$\Th{Q} \Proves
\eq[{\num k}' + \num n][\num m]$.
સમતાની સંક્રામકતા વડે
$\Th{Q} \Proves
\eq[{\num k}' + t_1][t_2]$;
એટલે $\Th{Q} \Proves
\lexists[z][\eq[z' + t_1][t_2]]$.
$!Q_8$ની જમણેથી ડાબી
દિશા પ્રમાણે $\Th{Q} \Proves
t_1 < t_2$.
\sourcecorrection{OLINC-033}{સ્થિર અંગ્રેજી
મૂળ પહેલા $\num n+\num k'$
નું સમીકરણ આપે છે, પરંતુ
પછીના અસ્તિત્વવાચક સાક્ષી માટે
$\num k'+\num n$ જરૂરી છે.
$\Th{Q}$ સામાન્ય સરવાળાની
અદલાબદલી સાબિત કરતું નથી;
અહીં નિશ્ચિત અંકોનું જરૂરી
સમીકરણ લેમ્માના પહેલા
ભાગથી સીધું લીધું છે.}

તેના બદલે $\Sat/{N}{t_1 < t_2}$,
એટલે $m \leq n$, માનો.
%
$\Th{Q}$માં રહીને
$t_1 < t_2$ ધારો. $!Q_8$ની
ડાબેથી જમણી દિશા પ્રમાણે
કોઈ~$z$ માટે
$\eq[z' + t_1][t_2]$.
$\Th{Q} \Proves
\eq[t_1][\num n]$ અને
$\Th{Q} \Proves \eq[t_2][\num m]$
હોવાથી
$\eq[z' + \num n][\num m]$.
%
હવે $!Q_5$ વાપરીને~$m$
પર બાહ્ય અનુમાન કરીએ, તો
$\eq[z' + \num{n - m}][\Obj 0]$.
જો $m = n$, તો $!Q_4$
વડે $\eq[z'][\Obj 0]$ મળે,
જે $!Q_2$ને વિરુદ્ધ છે.
જો $m < n$, તો ફરી~$!Q_5$
વડે
$\eq[(z' + \num{n - m - 1})'][\Obj 0]$,
જે $!Q_2$ને વિરુદ્ધ છે.
આથી $\Th{Q} \Proves
\lnot(t_1 < t_2)$.
\sourcecorrection{OLINC-034}{સ્થિર અંગ્રેજી
મૂળ આ બંને વિરોધાભાસને
$Q_3$ સાથે જોડે છે; ઉત્તરગામી
શૂન્ય ન હોય તે સ્વયંસિદ્ધિ
$Q_2$ છે. $m=n$ના
કિસ્સામાં સરવાળામાંથી શૂન્ય
દૂર કરવા $Q_4$ પણ જરૂરી છે.}
\end{proof}

\begin{lem}
\ollabel{lem:bounded-quant-equiv}
માનો કે $!A$ એ !!a{formula},
$t$ એ બંધ પદ અને
$k=\Value{t}{N}$. તો
\begin{enumerate}
\item $\Th{Q} \Proves \bforall{x<t}{!A(x)}$
ત્યારે અને માત્ર ત્યારે જ
$\Th{Q} \Proves
    !A(\num 0) \land \dots \land !A(\num{k-1})$.
\item $\Th{Q} \Proves \bexists{x<t}{!A(x)}$
ત્યારે અને માત્ર ત્યારે જ
$\Th{Q} \Proves
    !A(\num 0) \lor \dots \lor !A(\num{k-1})$.
\end{enumerate}
\end{lem}

\begin{proof}
    મર્યાદિત સર્વવાચક પરિમાણકનો
    કિસ્સો સાબિત કરીએ.
    $\Value{t}{N} = 0$ હોય,
    તો સમકક્ષતાની ડાબી બાજુ
    $\Th{Q}$માં સાબિત થાય છે,
    કારણ કે \olref[inc][req][min]{lem:less-zero}
    પ્રમાણે $x<\num 0$
    એવું શક્ય નથી.
    તેવી જ રીતે ખાલી સંયોજનને
    માત્ર $\ltrue$ લઈ શકીએ,
    જે પણ $\Th{Q}$માં
    સાબિત થાય છે.
    \sourcecorrection{OLINC-035}{સ્થિર અંગ્રેજી
    મૂળ આ સર્વવાચક કિસ્સામાં
    ખાલી વિયોજનને $\ltrue$
    કહે છે. અહીં સાન્ત
    સર્વવાચક દાવાને અનુરૂપ
    ખાલી સંયોજન જ $\ltrue$
    છે; ખાલી વિયોજનનો અર્થ
    જુદો થાય છે.}
    %
    તેથી કોઈ પ્રાકૃતિક સંખ્યા~$k$
    માટે $\Value{t}{N} = k+1$
    ધારો. \olref{lem:q-proves-clterm-id}
    વડે આપણે
    $\bforall{x<\num{k+1}}{!A(x)}$
    આકારના !!a{formula} સાથે
    કામ કરીએ તેમ માની શકીએ.
    
    માનો કે $\Th{Q} \Proves
    \bforall{x<\num{k+1}}{!A(x)}$,
    અને $n \leq k$ લો.
    \olref{lem:atomic-completeness}
    પ્રમાણે $\Th{Q} \Proves
    \num n < \num{k+1}$;
    તેથી તર્કથી $\Th{Q} \Proves
    !A(\num n)$. દરેક
    $n \leq k$ માટે આ વાત
    $k+1$ વાર લગાડવાથી
    માંગ્યા પ્રમાણે
    $\Th{Q} \Proves !A(\num 0)
    \land \dots \land !A(\num k)$.
    
    બીજી દિશા માટે
    $\Th{Q} \Proves
    !A(\num 0) \land \dots \land !A(\num k)$
    ધારો. $\Th{Q}$માં
    $x < \num{k+1}$ માનો.
    \olref[inc][req][min]{lem:less-nsucc}
    વડે
    $x = \num 0 \lor \dots \lor x = \num k$;
    તેથી તર્કથી~$!A(x)$
    મળે છે અને પરિણામે
    સર્વવાચક દાવો
    $\bforall{x<\num{k+1}}{!A(x)}$
    મળે છે.
    
    મર્યાદિત અસ્તિત્વવાચક
    પરિમાણિત !!{formula}sની
    સમકક્ષતાની સાબિતી પણ આવી જ છે.
\end{proof}

\begin{prob}
\olref{lem:bounded-quant-equiv}માં
અસ્તિત્વવાચક કિસ્સાની વિગતવાર
સાબિતી આપો.
\end{prob}

\begin{lem}
\ollabel{lem:delta0-completeness}
$!A$ એ $\Struct{N}$માં સાચું
$\Delta_0$ !!{sentence} હોય,
તો $\Th{Q} \Proves !A$.
\end{lem}

\begin{proof}
!!{formula}ની જટિલતા પર
અનુમાન વડે સાબિત કરીએ.
%
આરંભિક કિસ્સો
\olref{lem:atomic-completeness}
આપે છે; તેથી અનુમાનના
પગલા તરફ જઈએ. સરળતા માટે,
નિષેધ લાગુ પડે તે
!!{formula}ની રચના મુજબ
નિષેધના કિસ્સાને પેટાકિસ્સામાં
વહેંચીએ.

\begin{enumerate}
\item માનો કે $(!A \land !B)$
$\Struct{N}$માં સાચું છે;
ત્યારે $!A$ અને~$!B$ બંને
$\Struct{N}$માં સાચાં છે.
અનુમાનની ધારણા પ્રમાણે
$\Th{Q} \Proves !A$ અને
$\Th{Q} \Proves !B$;
તેથી તર્કથી
$\Th{Q} \Proves (!A \land !B)$.
%
\item માનો કે
$\lnot (!A \land !B)$ એ
$\Struct{N}$માં સાચું છે;
ત્યારે $\lnot !A$ અથવા
$\lnot !B$ એ $\Struct{N}$માં
સાચું છે. વ્યાપકતા ગુમાવ્યા
વિના પહેલું સાચું છે એમ
ધારો. અનુમાનની ધારણા પ્રમાણે
$\Th{Q} \Proves \lnot !A$;
તેથી તર્કથી $\Th{Q} \Proves
\lnot (!A \land !B)$.
%
\item માનો કે $(!A \lor !B)$
$\Struct{N}$માં સાચું છે;
ત્યારે $!A$ એ $\Struct{N}$માં
સાચું છે અથવા~$!B$ એ
$\Struct{N}$માં સાચું છે.
વ્યાપકતા ગુમાવ્યા વિના પહેલું
સાચું છે એમ ધારો.
અનુમાનની ધારણા પ્રમાણે
$\Th{Q} \Proves !A$;
તેથી તર્કથી $\Th{Q} \Proves
(!A \lor !B)$.
%
\item માનો કે
$\lnot(!A \lor !B)$ એ
$\Struct{N}$માં સાચું છે;
ત્યારે $\lnot !A$ અને
$\lnot !B$ બંને $\Struct{N}$માં
સાચાં છે. અનુમાનની ધારણા
પ્રમાણે $\Th{Q} \Proves
\lnot !A$ અને
$\Th{Q} \Proves \lnot !B$.
પરિણામે તર્કથી
$\Th{Q} \Proves
\lnot(!A \lor !B)$.
%
\item માનો કે
$\bforall{x<t}{!A(x)}$
$\Struct{N}$માં સાચું છે,
જ્યાં~$t$ બંધ પદ અને
$k=\Value{t}{N}$. અનુમાનની
ધારણા તથા તર્ક પ્રમાણે
દરેક $n < \Value{t}{N}$
માટે $!A(\num n)$ એ
$\Struct{N}$માં સાચું હોય,
તો $\Th{Q} \Proves
!A(\num 0) \land \dots \land !A(\num{k-1})$.
\olref{lem:bounded-quant-equiv}
પ્રમાણે $\Th{Q} \Proves
\bforall{x<t}{!A(x)}$.
%
\item મર્યાદિત અસ્તિત્વવાચક
પરિમાણકનો કિસ્સો, જેમાં
$\bexists{x < t}{!A(x)}$
આકારનું !!a{sentence} હોય,
તે મર્યાદિત સર્વવાચક
પરિમાણકના કિસ્સા જેવો જ છે.
%
\item માનો કે
$\lnot \bforall{x<t}{!A(x)}$
$\Struct{N}$માં સાચું છે,
જ્યાં~$t$ બંધ પદ છે.
આ !!{sentence} એ
!!{sentence}
$\bexists{x<t}{\lnot !A(x)}$
ની સમકક્ષ છે, અને સમકક્ષતા
$\Th{Q}$માં નિગમ્ય છે;
તેથી મર્યાદિત અસ્તિત્વવાચક
પરિમાણકની દલીલ વાપરી શકીએ.
%
\item તેવી જ રીતે માનો કે
$\lnot \bexists{x<t}{!A(x)}$
$\Struct{N}$માં સાચું છે,
જ્યાં~$t$ બંધ પદ છે.
આ !!{sentence} એ
$\Th{Q}$માં
$\bforall{x<t}{\lnot!A(x)}$
ની સમકક્ષ છે; તેથી
મર્યાદિત સર્વવાચક પરિમાણકની
દલીલ વાપરી શકીએ.
\sourcecorrection{OLINC-036}{સ્થિર અંગ્રેજી
મૂળના આ કિસ્સામાં
$\bexists$ માટે બીજો
કૌંસબંધ દલીલ-ઘટક છૂટી
ગયો છે. મેક્રોની વ્યાખ્યા
પ્રમાણે અહીં $!A(x)$ને
બીજા ઘટકમાં મૂક્યું છે.}
%
\item છેલ્લે માનો કે
$\lnot !A$ એ $\Struct{N}$માં
સાચું છે. હવે માત્ર બે
કિસ્સા બાકી છે: $!A$
પરમાણુ છે, અથવા કોઈ
$\Delta_0$ !!{sentence}~$!B$
માટે
$\lnot !A \ident \lnot\lnot !B$.
જો $!A$ પરમાણુ હોય, તો
\olref{lem:atomic-completeness}
પ્રમાણે $\Th{Q} \Proves
\lnot !A$. જો
$\lnot !A \ident \lnot\lnot !B$,
તો તર્કથી તે~$!B$ની
$\Th{Q}$માં સાબિત કરી શકાય
એવી સમકક્ષતા ધરાવે છે.
તે $\Struct{N}$માં સાચું છે,
કારણ કે $\lnot !A$ પણ
$\Struct{N}$માં સાચું છે.
તેથી અનુમાનની ધારણા પ્રમાણે
$\Th{Q} \Proves \lnot !A$.
\end{enumerate}
\end{proof}

\begin{prob}
\olref{lem:delta0-completeness}માં
અસ્તિત્વવાચક કિસ્સાની વિગતવાર
સાબિતી આપો.
\end{prob}

\begin{thm}
\ollabel{thm:sigma1-completeness}
$!A$ એ $\Struct{N}$માં સાચું
$\Sigma_1$ !!{sentence} હોય,
તો $\Th{Q} \Proves !A$.
\end{thm}

\begin{proof}
$\lexists[x][!A(x)]$ એ
$\Struct{N}$માં સાચું
$\Sigma_1$ !!{sentence}
હોય, તો કોઈ પ્રાકૃતિક
સંખ્યા~$n$ અને ચલ-નિયુક્તિ~$s$
એવાં છે કે $s(x) = n$ અને
$\Sat{N}{!A(x)}[s]$.
સંતોષસંબંધનાં પ્રમાણભૂત
તથ્યો વડે $\Sat{N}{!A(\num n)}$
મળે છે. પરંતુ $!A(\num n)$
$\Delta_0$ !!{formula} છે;
તેથી \olref{lem:delta0-completeness}
પ્રમાણે $\Th{Q} \Proves
!A(\num n)$ અને પરિણામે
તર્કથી $\Th{Q} \Proves
\lexists[x][!A(x)]$ પણ છે.
\sourcecorrection{OLINC-037}{સ્થિર અંગ્રેજી
મૂળ અહીં $\lexists{x}!A(x)$
લખે છે; આ મેક્રો ચલ અને
વ્યાપ માટે ચોરસ કૌંસવાળા
વૈકલ્પિક ઘટકો લે છે.
વ્યાખ્યા અને સાબિતીના
છેલ્લા સૂત્ર પ્રમાણે
$\lexists[x][!A(x)]$
આકાર સ્પષ્ટ કર્યો છે.}
\end{proof}

\end{document}
